\chapter*{Summary}
Recommender systems (RecSys) have been widely deployed in online platforms to help users discover relevant items from large collections. Despite their success and the tremendous business value they generate, existing RecSys still face several challenges, including data sparsity and the limited ability to incorporate rich semantic knowledge beyond user–item interactions. This dissertation studies how to improve personalized recommendation by integrating structural learning from interaction graphs with semantic knowledge derived from large language models (LLMs). Specifically, five research topics are investigated. First, a unified hypergraph pretraining framework is proposed to incorporate heterogeneous relational signals into recommendation models through multitask learning. Second, textual prompts encoded by pretrained language models are introduced to guide hypergraph pretraining for recommendation. Third, a token-level alignment approach is developed to transfer semantic knowledge from LLM representations to graph-based recommendation models. Fourth, LLMs are used to infer substitute and complementary relations between items to refine relational signals in recommendation. Finally, an LLM-based agent framework is proposed to integrate multiple recommendation tools and perform reasoning over item relations during recommendation. Experiments are conducted on several public benchmark datasets including Amazon, Steam, Goodreads and Instacart datasets. The experiments involve only computational analysis of existing datasets and do not involve human subjects, animals, or recombinant DNA. Therefore, approvals from the Institutional Review Board, the Animal Care Committee, or the Institutional Biosafety Committee are not required.